% Generated draft: requires mathematical and source review.
\documentclass[11pt]{article}
\usepackage[margin=25mm]{geometry}
\usepackage{fontspec}
\setmainfont{Latin Modern Roman}
\usepackage{amsmath,amssymb,mathtools,unicode-math}
\setmathfont{Latin Modern Math}
\usepackage{enumitem,needspace,xurl,hyperref}
\setlength{\parindent}{0pt}
\setlength{\parskip}{5pt}
\setlength{\emergencystretch}{3em}
\newcommand{\N}{\mathbb{N}}
\newcommand{\Z}{\mathbb{Z}}
\newcommand{\Q}{\mathbb{Q}}
\newcommand{\R}{\mathbb{R}}
\newcommand{\C}{\mathbb{C}}
\newcommand{\F}{\mathbb{F}}
\newcommand{\A}{\mathbb{A}}
\newcommand{\PP}{\mathbb{P}}
\newcommand{\E}{\mathbb{E}}
\newcommand{\eps}{\varepsilon}
\newcommand{\dd}{\,\mathrm{d}}
\begin{document}
\section*{UnsolvedMath ID: 8800001}
\textbf{Source problem code:} catalogue entry
\subsection{Definitions and mathematical statement}
Let \(q\) be a prime power and \(E/\F_q\) an elliptic curve, given in a specified finite-field representation, with \(P\in E(\F_q)\). An isogeny \(\psi:E\to E'\) is a nonconstant algebraic group morphism and has positive finite degree \(D=\deg\psi\). If \(E'\) and \(\psi\) are defined over \(\F_q\), the output \(\psi(P)\) lies in \(E'(\F_q)\); otherwise the output extension field and its representation must be included in the problem data.

The question is to characterize representations/families of such maps permitting uniform evaluation in time polynomial in
\[
\log D\quad\text{and}\quad\log q,
\]
ideally \(O(\log D\,(\log q)^2)\). The logarithmic-degree expression is an asymptotic target for \(D\ge2\); degree-one maps require a separate convention such as \(1+\log D\), not a literal zero-time bound. Constant maps have degree zero and are not included in this isogeny formulation.

The representation of \(\psi\) is essential. Dense rational functions of degree comparable to \(D\) can have description length far exceeding \(\operatorname{poly}(\log D,\log q)\); a succinct composition/circuit description is different input. The source gives neither encoding nor a cost for finding that encoding. It also does not separate preprocessing of \(\psi\) from online evaluation at \(P\). Unrestricted precomputed tables would make a small online lookup time an inadequate characterization. A precise theorem must therefore account for input length, preprocessing, field-operation versus bit complexity, and uniformity over \(P\).

The source also says maps or isogenies, without specifying the larger map class, separability, or field of definition. Those are not silently restricted to ordinary curves or to a favorable endomorphism case. The supplied research summary reports both positive and negative results and says no general logarithmic-degree algorithm/characterization is established, especially in large-discriminant ordinary settings; it gives no precise theorem or lower-bound reference.
\subsection{Short English statement}
Which succinctly represented elliptic-curve isogenies can be evaluated in time polynomial in \(\log\deg\psi\) and \(\log q\), ideally near the stated product bound? Encoding, preprocessing, field of definition, and computation model are part of the question.
\subsection{Sources}
\begin{itemize}
\item Local UnsolvedMath registry, permanent ID \(8800001\), exact source code \texttt{catalogue entry}, ``Fast evaluation of high-degree elliptic-curve isogenies.'' Proposer Denis Charles.
\item IPAM \emph{Workshop on Number Theory and Cryptography --- Open Problems} (2006), attachment 6218, recorded PDF page 1; schedule source \url{https://www.ipam.ucla.edu/programs/workshops/workshop-i-number-theory-and-cryptography-open-problems/?tab=schedule}. Recorded access 2026-07-29; pdf-and-legacy-ppt visual-audit extraction.
\item Supplied PARTIAL-PROGRESS research summary states positive/negative results exist but supplies no exact algorithm, lower bound, or complete publication citation.
\end{itemize}
\end{document}
